\chapter{Ажлын төлөвлөгөө}
ШМТК-ийн Компьютерын ухааны тэнхимд сурацаж буй Онолбаатар овогтой Дөлгөөний "\titlename" сэдэвт 
судалгааны ажлыг хийж гүйцэтгэх төлөвлөгөө.

\setlength{\arrayrulewidth}{0.4pt}
\begin{longtable}{|p{2.2cm}|p{2cm}|p{10.5cm}|}
\hline
\textbf{Сар} & \textbf{Долоо хоног} & \textbf{Төлөвлөгөө} \\
\hline
\endfirsthead
\hline
\textbf{Сар} & \textbf{Долоо хоног} & \textbf{Төлөвлөгөө} \\
\hline
\endhead
\hline
\endfoot

\multirow{4}{*}{10 сар} & 1 & Сэдвийн анхны сонголт, ажлын хүрээ тодорхойлох. \\
\cline{2-3}
& 2 & Сэдвийн хүрээ, судалгааны асуултуудыг тодорхойлох. \\
\cline{2-3}
& 3 & Онол, арга зүйн материал судлах; дүрс/нүүр таних алгоритмуудтай танилцах. \\
\cline{2-3}
& 4 & OpenCV, MTCNN сан болон алгоритмуудын туршилт хийх. \\
\hline

\multirow{5}{*}{11 сар} & 1 & Python орчинд камерын өгөгдөл боловсруулах туршилт. \\
\cline{2-3}
& 2 & Нүүр илрүүлэх ба таних загвар боловсруулах; өгөгдлийн сан үүсгэх. \\
\cline{2-3}
& 3 & Алгоритмын үр дүнгийн үнэлгээ. \\
\cline{2-3}
& 4 & Системийн туршилтын анхны хувилбар. \\
\cline{2-3}
&  & \textcolor{red}{TCA үзлэг 1}. \\
\hline

\multirow{4}{*}{12 сар} & 1 & Туршилтын өгөгдөл нэмэх; системийн сайжруулалт. \\
\cline{2-3}
& 2 & Ирц бүртгэлийн өгөгдлийн сан болон интерфэйс боловсруулах. \\
\cline{2-3}
& 3 & Алгоритмын сайжруулалт; анхны дүнгийн үнэлгээ. \\
\cline{2-3}
& 4 & Өгөгдлийн сан цэвэрлэгээ, давхардал арилгах, анхны gallery rebuild хийх. \\
\hline

\multirow{4}{*}{1 сар} & 1 & Алгоритмын сайжруулалт; камер + өгөгдлийн сангийн нэгдсэн туршилт. \\
\cline{2-3}
& 2 & Алгоритмын сайжруулалт; камер + өгөгдлийн сангийн нэгдсэн туршилт. \\
\cline{2-3}
& 3 & Threshold tuning, validation sweep, туршилтын үр дүнгийн анхны нэгтгэл. \\
\cline{2-3}
& 4 & Тайлан бичих; дүгнэлт боловсруулах. \\
\hline

\multirow{4}{*}{2 сар} & 1 & Тайлан сайжруулалт. \\
\cline{2-3}
& 2 & Тайлан сайжруулалт; алдааны шинжилгээ, дэлгэцийн зураг ба тайлбарт тэмдэглэгээ бэлтгэх. \\
\cline{2-3}
& 3 & Туршилтын үр дүнг нэгтгэх; алгоритмын эцсийн сайжруулалт. \\
\cline{2-3}
& 4 & Туршилтын үр дүнг нэгтгэх; алгоритмын эцсийн сайжруулалт. \\
\hline

\multirow{4}{*}{3 сар} & 1 & \textcolor{red}{TCA үзлэг 2}. \\
\cline{2-3}
& 2 & Дипломын бүлгүүдийн найруулга сайжруулах; ном зүй, форматын засвар хийх. \\
\cline{2-3}
& 3 & Ёс зүй, хувийн мэдээлэл хамгаалалтын хэсэг нэмэх; дүгнэлтийн ноорог бэлтгэх. \\
\cline{2-3}
& 4 & Эцсийн валидацын шалгалт хийх; зураг, хүснэгтийн шинэчлэл хийх. \\
\hline

\multirow{4}{*}{4 сар} & 1 & Эцсийн засвар, эхний draft-ын алдааг цэгцлэх. \\
\cline{2-3}
& 2 & Хураангуйг шинэчлэн бичих; бодит туршилтын дүнгээр үр дүнгийн бүлгийг шинэчлэх. \\
\cline{2-3}
& 3 & Хамгаалалтын илтгэлийн слайд болон үзүүлэн материал бэлтгэх. \\
\cline{2-3}
& 4 & \textcolor{red}{TCA урьдчилсан хамгаалалт}. \\
\hline

\multirow{4}{*}{5 сар} & 1 & Эцсийн эх бэлдэх, хэвлэхэд бэлтгэх. \\
\cline{2-3}
& 2 & Эцсийн үг үсгийн хяналт хийх; хэвлэлийн бэлтгэл хангах. \\
\cline{2-3}
& 3 & Хамгаалалтын урьдчилсан бэлтгэл хийх; эцсийн засвар оруулах. \\
\cline{2-3}
& 4 & \textcolor{red}{TCA үндсэн хамгаалалт}. \\
\hline
\end{longtable}


